\section{Mathematical Background}

\subsection{Computably Enumerable Sets}

A computably enumerable set can be presented by an effective monotone
sequence of finite approximations
\[
  A_0 \subseteq A_1 \subseteq A_2 \subseteq \cdots
\]
and the limit set \(A=\bigcup_s A_s\).  The monotonicity is positive:
membership can change from absent to present when an element is enumerated,
but never from present to absent.

This distinction matters for formalization.  A prefix-monotone Boolean
string representation would fix every bit once it appears in the prefix.
Such a representation is appropriate for some finite-information
arguments, but not for c.e. approximations themselves: a c.e. construction
must be allowed to discover later that a previously absent number should be
enumerated.  In the Lean development, stage approximations are finite sets
derived from lists, while Boolean lists are used only as bounded oracle
snapshots for individual computations.

\subsection{Oracle Computation and Use}

For an oracle set \(B\), the notation \(\Phi_e^B(x)\) denotes the result of
running the \(e\)-th oracle program on input \(x\) with access to \(B\).  A
halting computation uses only finitely many oracle answers.  Its
\emph{use} is an upper bound on the queried positions.  The formalization
uses the strict convention: if position \(p\) is queried, then the recorded
use is at least \(p+1\), and query-free computations have use \(0\).

The use principle is the basis of restraint.  If a computation halts with
use \(u\), then any later oracle agreeing below \(u\) produces the same
finite computation and the same output.  A priority strategy can therefore
preserve a computation by preventing future enumeration below \(u\).

\subsection{The Friedberg--Muchnik Requirements}

The construction builds two c.e. sets \(A\) and \(B\).  The requirements are
interleaved by natural priority index:
\[
  R_e : A \neq \Phi_e^B,\qquad
  S_e : B \neq \Phi_e^A,
\]
with \(R_e\) represented by priority \(2e\) and \(S_e\) by priority
\(2e+1\).  Lower numerical priority indices have higher priority.

The strategy for \(R_e\) chooses a fresh witness \(x\), keeps \(x\notin A\),
and waits for a stage at which the current finite approximation to
\(\Phi_e^B(x)\) converges to \(0\).  If that happens, it enumerates \(x\)
into \(A\), records the computation's use as a restraint on \(B\), and
initializes all lower-priority requirements.  The \(S_e\) strategy is the
mirror image, exchanging \(A\) and \(B\).

\subsection{Finite Injury}

The classical finite-injury proof proceeds by induction over priority.
Once all higher-priority requirements stop acting, the current requirement
can no longer be initialized from above.  It then makes only finitely many
progress steps: appoint a witness, possibly act, and then remain stable.
The Lean proof makes this argument precise using the definitions
\texttt{rank} and \texttt{reqRank}; after higher-priority requirements are
quiet, each attention to the current requirement strictly increases a rank
bounded by \(2\).

